% from ega4-2-fr.tex
\oldpage[IV-2]{98}

Soient $\mathfrak p$, $\mathfrak p'$ les images réciproques de
$\mathfrak q$, $\mathfrak q'$ respectivement dans $A$, $\mathfrak n$ le
noyau de l'homomorphisme $A_{\mathfrak p}\to B_{\mathfrak q}$.

L'image $A'$ de $A_{\mathfrak p}$ dans $B_{\mathfrak q}$ étant isomorphe à
$A_{\mathfrak p}/\mathfrak n$, la formule (5.6.1.1) appliquée à $A'$ et à
$B'=B_{\mathfrak q}$ s'écrit
\begin{equation}
\dim(A_{\mathfrak p}/\mathfrak n)
+\operatorname{deg.tr.}_{A_{\mathfrak p}/\mathfrak n}(B_{\mathfrak q})
=\dim(B_{\mathfrak q})
+\operatorname{deg.tr.}_{k(\mathfrak p)}k(\mathfrak q).
\tag{5.6.1.3}\label{IV.2.s5.6.1.3-fr}
\end{equation}

D'autre part, le noyau de $A_{\mathfrak p'}\to B_{\mathfrak q'}$ est
$\mathfrak nA_{\mathfrak p'}$, donc, en appliquant la formule (5.6.1.1) où
l'on remplace $A'$ par $A_{\mathfrak p'}/\mathfrak nA_{\mathfrak p'}$ et
$B'$ par $B_{\mathfrak q'}$, on a
\begin{equation}
\dim(A_{\mathfrak p'}/\mathfrak nA_{\mathfrak p'})
+\operatorname{deg.tr.}_{A_{\mathfrak p'}/\mathfrak nA_{\mathfrak p'}}
(B_{\mathfrak q'})
=\dim(B_{\mathfrak q'})
+\operatorname{deg.tr.}_{k(\mathfrak p')}k(\mathfrak q').
\tag{5.6.1.4}\label{IV.2.s5.6.1.4-fr}
\end{equation}

Enfin, $B/\mathfrak q'$ est une $A$-algèbre intègre de type fini, l'image
réciproque de $\mathfrak q/\mathfrak q'$ dans $A$ est $\mathfrak p$, et le
noyau de l'homomorphisme
$A_{\mathfrak p}\to B_{\mathfrak q}/\mathfrak q'B_{\mathfrak q}$ est
$\mathfrak p'A_{\mathfrak p}$, donc, en appliquant (5.6.1.1) où l'on
remplace $A'$ par $A_{\mathfrak p}/\mathfrak p'A_{\mathfrak p}$ et $B'$ par
$B_{\mathfrak q}/\mathfrak q'B_{\mathfrak q}$, on a
\begin{equation}
\dim(A_{\mathfrak p}/\mathfrak p'A_{\mathfrak p})
+\operatorname{deg.tr.}_{A_{\mathfrak p}/\mathfrak p'A_{\mathfrak p}}
(B_{\mathfrak q}/\mathfrak q'B_{\mathfrak q})
=\dim(B_{\mathfrak q}/\mathfrak q'B_{\mathfrak q})
+\operatorname{deg.tr.}_{k(\mathfrak p')}k(\mathfrak q).
\tag{5.6.1.5}\label{IV.2.s5.6.1.5-fr}
\end{equation}

Ajoutons maintenant (5.6.1.4) et (5.6.1.5) membre à membre, et notons que
$k(\mathfrak p')$ (resp. $k(\mathfrak q')$) est le corps des fractions de
$A_{\mathfrak p}/\mathfrak p'A_{\mathfrak p}$ (resp.
$B_{\mathfrak q}/\mathfrak q'B_{\mathfrak q}$); par ailleurs
$A_{\mathfrak p'}/\mathfrak nA_{\mathfrak p'}$ et
$A_{\mathfrak p}/\mathfrak n$ ont même corps des fractions, et
$B_{\mathfrak q'}$ et $B_{\mathfrak q}$ ont même corps des fractions;
enfin, puisque $A$ est supposé caténaire, on a (0, 16.1.4.2)
\[
\dim(A_{\mathfrak p}/\mathfrak p'A_{\mathfrak p})
+\dim(A_{\mathfrak p'}/\mathfrak nA_{\mathfrak p'})
=\dim(A_{\mathfrak p}/\mathfrak n).
\]
Tenant compte de ces remarques et de (5.6.1.3), il vient la relation
(5.6.1.2). C.Q.F.D.

\begin{definition}[5.6.2]
\label{IV.2.s5.6.2-fr}
On dit qu'un anneau noethérien $A$ est \emph{universellement caténaire} s'il
vérifie les conditions équivalentes a), b), c) de (5.6.1).
\end{definition}

\begin{remarks}[5.6.3]
\label{IV.2.s5.6.3-fr}
\begin{enumerate}
\item[(i)] Si $A$ est universellement caténaire, il en est de même de
$S^{-1}A$ pour toute partie multiplicative $S$ de $A$, comme il résulte
aussitôt de (5.6.1, a)) et du fait que tout anneau de fractions d'un anneau
caténaire est caténaire. Inversement, si, pour tout idéal premier
$\mathfrak p$ de $A$, l'anneau $A_{\mathfrak p}$ est universellement
caténaire, alors $A$ est universellement caténaire: cela résulte de
(5.6.1, c)), si l'on remarque qu'en posant $S=A-\mathfrak p$, $S^{-1}B$
est une $A_{\mathfrak p}$-algèbre de type fini, et
$B_{\mathfrak q}=(S^{-1}B)_{S^{-1}\mathfrak q}$.
\item[(ii)] On dit qu'un préschéma localement noethérien $X$ est
\emph{universellement caténaire} si, pour tout $x\in X$, l'anneau
$\mathcal O_x$ est universellement caténaire. Il résulte de (i) que pour que
$A$ soit un anneau universellement caténaire, il faut et il suffit que le
schéma $\operatorname{Spec}(A)$ le soit.
\item[(iii)] Le critère (5.6.1, c)) ne fait intervenir que les images de $A$
dans des $A$-algèbres intègres de type fini, donc des anneaux quotients
intègres de $A$. On en conclut que si

